\subsection{Scalar-bearing trace estimates with the prescribed normalizer}
\label{sec:retained-scalar-weighted-traces}

The remaining common-family comparison must retain the scalar
moduli in Section~\ref{sec:retained-column-positive-traces}.
They need not, however, be estimated by a uniform maximum
over all words. The following two elementary estimates
preserve the prescribed normalizer and allow a weighted
arithmetic bound.

Let \(T_i:V_i\to W\) be maps of finite-dimensional Hermitian
spaces, with \(d_i=\dim V_i\), \(R=\sum_i d_i>0\),
\(D=\dim W>0\), \(c>0\). Write
\(S_i=T_i^*T_i\), \(Q_i=T_iT_i^*\),
\(m_i=\operatorname{tr}S_i=\operatorname{tr}Q_i\).
For scalar factors \(a_i\) put \(\rho_i=|a_i|^2\),
\(S_\rho=\bigoplus_i\rho_i S_i\), \(Q_\rho=\sum_i\rho_i Q_i\).
For any such family let \(g,b\) be the normalized positive
source deficit and target excess, both with the specified \(c\).

\begin{proposition}[Weighted scalar stability]
\label{prop:retained-weighted-scalar-stability}
For the rescaled maps \(a_iT_i\) one has
\[
 \begin{split}
 c g_\rho&\le cg+\sum_i(1-\rho_i)_+m_i,\\
 c b_\rho&\le cb+\sum_i(\rho_i-1)_+m_i,\\
 g_\rho+b_\rho&\le g+b+c^{-1}\sum_i|\rho_i-1|m_i .
 \end{split}                                                       \tag{ST1}
\]
These inequalities hold as stated after any prescribed
source and target orthogonal cuts, with \(m_i\) the
mass of the actual compressed maps and \(R,D,c\) their
prescribed dimensions and normalizer.
\end{proposition}
\begin{proof}
The function \(A\mapsto\operatorname{tr}A_+\) on Hermitian
matrices is monotone in the positive-semidefinite order
and subadditive, by its variational formula
\(\sup_{0\le P\le1}\operatorname{tr}(PA)\).
The source difference is bounded above by
\(\bigoplus_i(1-\rho_i)_+S_i\), and the target difference
by \(\sum_i(\rho_i-1)_+Q_i\).
Apply these facts to the two positive parts. Scalars
commute with the prescribed cuts, so the same argument
works for compressed maps. Adding the two allowances
charges each absolute scalar deviation exactly once.
\end{proof}

In particular suppose the exact common-family comparison
on the retained spaces gives \(T_i^{\rm actual}=a_iT_i\)
up to source unitaries and a common target isometry.
If the unscaled family has proved budget \(E\), then
the contribution to the final norm penalty is at most
\[
 \kappa^{-1}\operatorname{avg}_u
 \min\left\{H_u,\,
  \left(E_u+c_u^{-1}\sum_i|\rho_{ui}-1|m_{ui}\right)
                      (U_u+H_u)\right\}.                  \tag{ST2}
\]
The scalar-deviation allowance may itself be bounded in this
clipped weighted average. Neither an individual
\(\rho_{ui}=1\) nor a uniform maximum of their deviations
is necessary. A choice of a common alternative reference
scale is harmless only if its positive traces are computed
relative to the actual \(c_u\).

There is also a direct weighted second-moment formulation,
useful when the scalars are not close to one. Suppose
\[
 \|S_i-1_{V_i}\|_{\rm HS}\le e_i\sqrt{d_i},\qquad
 \|T_i^*T_j\|_{\rm HS}^2\le
          (1+\eta)d_i d_j/D_*\quad(i\ne j)
\]
for the actual comparison dimension \(D_*>0\).
Set
\[
 M_0=\sum_i\rho_i d_i,\quad
 \Delta=\sum_i\rho_i e_i d_i,\quad
 B_\rho=\sum_i\rho_i^2(1+e_i)^2d_i,\quad
 C_\rho=\sum_i\rho_i^2d_i^2 .
\]
Then the following fully explicit bound holds:
\begin{equation}\label{eq:retained-weighted-second-moment}
 g_\rho+b_\rho\le\frac{1}{2c}
 \left[
 \sum_i|\rho_i-c/R|d_i+\Delta+
 \sqrt{D B_\rho+
   \frac{D(1+\eta)}{D_*}(M_0^2-C_\rho)-2c(M_0-\Delta)+c^2}
 \right].
\end{equation}
Indeed the source trace norm is at most the sum inside
the first two terms by the triangle inequality and the
trace/Hilbert--Schmidt inequality on each block.
The actual common mass is in
\([M_0-\Delta,M_0+\Delta]\).
For the target,
\(\operatorname{tr}Q_\rho^2=
\sum_{i,j}\rho_i\rho_j\|T_i^*T_j\|_{\rm HS}^2\);
the diagonal terms give \(B_\rho\), the off-diagonal
terms give \((1+\eta)(M_0^2-C_\rho)/D_*\).
Now use the exact prescribed-normalizer trace-norm identity
(M1) of the preceding section, bounding the target trace
norm by the square root in its actual dimension \(D\).
The radicand is nonnegative, since it majorizes the
square of the target Hilbert--Schmidt error times \(D\).
No row-wise maximum of \(e_i\) or \(\rho_i\) occurs.

These results are sufficient trace and budget reductions,
not a construction of the actual \(a_i\), the common target,
the reference dimension, or the off-diagonal estimates.
The same-data anchored see-saw and moving concentration
hypotheses remain required. In particular none of these
inequalities justifies using a different \(c\) in an
arithmetic height or covolume comparison.
